\section{Hermetism and Esoteric Christianity Rise to Converge}

Gornahoor has conducted a number of thought exercises, designed to aid the aspirant in ``mental fasting'', a discipline of dropping old thoughts and thought-patterns or memes (compulsions). These exercises involve recovering a sense of Self that is already there, or (as it was put), laying aside the chain-mail of discursive thought, which often (merely) functions to jam any impulses coming from higher Being or higher Self to the ordinary, day-to-day self. It is argued that this is not merely an ``eastern'' practice but also a hermetic one. Perhaps, if the East is the land of dawn, the West of manifesting, it is the North that is origin, or ab-origin, and it is to Hyperborea that thought-meditation ought to tend, rather than to this or that guru, or this or that order.

Is Gornahoor off-base here? Let us test this --- what do we find when we examine a random twentieth century hermeticist? Let us take Franz Bardon, or ``Frabato''. Our source will be this recommended guide\footnote{\url{http://www.abardoncompanion.com/IIH-Intro\&Theory.html}} to interpreting his few works which he left. This is Rawn's Commentary on \emph{Initiation into Hermetics}.

\begin{quotationx}
In the initial exercises of Step One, Bardon describes three sorts of mental discipline or meditation. The first type involves merely observing what goes on in your mind. In this exercise, the student does not block any thoughts, s/he merely observes what is present. Given time and repeated practice, you will notice that the flow of thoughts naturally slows down. But what is really happening is that you are re-tuning your mind to another, less cluttered, level of mentation. This is not something that you can force, so it does little good at this stage to be blocking certain thoughts while letting others through, etc.

Of concern here, is the other distractions that arise, such as that car alarm that keeps going off in the distance, or the bark of the neighbor's dog. These sorts of incidents can distract your attention from the observance of your thoughts. While such occurrences are not within your ability to control, your response to them is within your control. So, you must learn how to quickly dismiss these distractions and refocus your attention to the task at hand. At first this may be difficult, but with persistent practice, your ability to refocus will become so quick and absolute that such external events will no longer distract; or rather, the distraction will be so brief that it will not interrupt your practice.

Another sort of distraction is that you will be tempted to pursue the thoughts that arise in your mind. The point here however, is to distance yourself from involvement with your individual thoughts --- you are to be only an observer, not a participator. At first, this is also very difficult, but with persistent practice, you will learn how to distance yourself and observe.
\end{quotationx}

This is virtually identically to what Gornahoor has claimed in regards to ``enlightenment''. Gornahoor has attempted to direct attention away from ``the East'' as a kind of reflex knee-jerk reaction to Western fatigue. Bowdlerized Eastern discipline has not produced what was promised on our spiritual soil.

But what of Gornahoor's claims that esoteric Christianity has a real existence in the West, and is (in fact) the ``missing link'' between all of the variegated Western traditions? That, beginning with Clement and St. John and the Baptist and others, the West has not lacked for initiation and potent spiritual possibilities?

Let's examine another random Hermeticist, in this case, Tomasso Palamadessi\footnote{\url{http://www.archeosofica.org/en/content/view/152/53/}}, who founded the Archeosophical Society in Turin in 1948 to perpetuate and unite all of the disparate streams of influence in the West:

\begin{quotex}
The same human being when he has the proof of his own survival, with all the implied consequences, then he will want to defend at all costs this survival in the best possible conditions, not as a wandering and confused spirit, now in a fleshly body in this world and now in the immense cosmic ocean of the Hereafter, but as a personal, deified I, which at the end of its days can depart safe and tranquil, not as an astronaut, but as a Theonaut, to land in the Kingdom of Heaven where God has been waiting for him for always. This journey requires two interventions: God from on high, us from below: divine Grace and human ascesis.\footnote{Source, Pamphlet \#1\footnote{\url{http://www.archeosofica.org/en/content/view/152/53/}}}
\end{quotex}

Here we find the same emphasis on deification that Evola has made, but this time in a context that renders Evola more understandable.

Was there or was there not a valid secret oral succession from the earliest days, as Gornahoor has claimed, within the Faith?

\begin{quotex}
And what else could have been meant by men like Saint Basil, Father of the Greek Church, who lived in close contact with the initiated Monks of the Orient, in the year 374 AD in Caesarea of Cappadocia, today Kaisarije in Anatolia, by saying: ``We receive the dogmas that have been transmitted to us by right and those that have come to us from the Apostles under the veil and under the mystery of an oral tradition. How could be diffused publicly that which noninitiates are forbidden to contemplate? ... This is exactly why many things were transmitted in an unwritten form ... '' (Treatise on the Holy Spirit, XVII). And Clement in Stromata I, chapter XII specifies: ``Since the holy tradition could not be something common and public, at least if one realizes the greatness of its teachings, it is necessary to hide `this wisdom expressed in mystery', that the Son of God has taught us.''
\end{quotex}


It goes even further back, as suggested on Meditations on the Tarot.

\begin{quotex}
And with us there is always John, the Prophet of the Apocalypse, to whom Jesus left the task of developing the esoteric and invisible Ekklesia, discontinuous and ardent like the fire, instituted on the Calvary when John was made son of the Virgin Mary and brother of Christ, and confirmed on the lake of Tiberias. To Peter was entrusted the flock and the martyrs (lambs), but to John the more evolved humanity with the command of remaining until his return (John, XXI, 4-25).
\end{quotex}


This is astounding. Between the two of them, Bardon \& Palamadessi, Gornahoor's claims are not mere validated, but practically duplicate. How can this be?

The closer one gets to ``true North'', to origin, the less disagreement or disjunction can possibly remain, even in the deep recesses of the mind.

If one can lay aside the fact that he looks too much like Anthony Hopkins, Bardon (as interpreted here in this companion work) is recommending the ``royal path'' of alchemy, which is essentially a secular and individualized version of esoteric Christianity.

To whet the appetite, consider step \#4:


\begin{quotex}
The fourth and most complete stage of transplantation of consciousness can take years of work to master and is attainable only if one has strictly adhered to the highest moral code. With this fourth type, you switch from mere observer, to active participant. Not only do you experience the whole being (physical, astral and mental), but you also become one with the whole being. In effect, your mental body joins the mental body of the being you are transferring your consciousness into, and you thereby have access to an absolute degree of direct influence over the actions, emotions and thoughts of the being. However, the being into which you have transferred your consciousness to this degree, will also have access to your being from the inside out. Usually, the being you would transfer into in this manner, would not have the magical ability to even perceive your presence and take advantage of the connection you have established. Nonetheless, the consequences of the mutuality of such a connection should be considered beforehand.
\end{quotex}


This is (by Jove himself) precisely what Mouravieff argued needed to occur among the second-born, if a world holocaust in the Kali Yuga was to be avoided.

Shall we continue? Is it even necessary? Skim the literature, and more and more striking similarities of teaching, ideas, and practice will leap out.

Palamedessi is no less shocking --- he claims that the initiation of the pagan world were powerful (but real) illusions dependent upon the personality and force of the initiator, who created upon the astral body of the initiate (during a three day sleep) a vision of the other worlds. So that the difference between Christ and paganism was simply that Christ's initiation opens the possibility of actual or entirely potentiated initiation that does not in fact depend upon the person of Christ, after the period of tutelage (self-growth, rather than an Eleusian mystery sleep) is ended. We worship Christ, in order to become a little brother, who is a joint heir.

This does not denigrate the pagan mysteries --- on the contrary, it elevates them. They are, and remain, valid. They can be preparations for the real one, if the pagan has eyes to be born again, after being born again. Likewise, the Christian is given ``content'' to what could threaten to be an abstract faith. Christ has many masks.

\begin{quotationx}
In the Sanctuaries of India, Egypt, Greece, the disciple was put by his instructor into a state of lethargic sleep for three days and three nights, after a preliminary physical and moral catharsis (or purification), fasts and prayers. During the lethargy, the soul (which is composed of the three principles: spirit, pathetic soul and eros, plus the subtle bodies, namely the causal, mental, emotional and etheric) separated itself from the physical body for a certain time, though remaining connected to the body by means of the tenuous silver cord (see in the Bible Ecclesiastes, XII, 6).

Outside of the body the double, that is, the whole of the interpenetrating subtle bodies, being more impressionable and moldable, was for the Initiator the matter on which he could act with his mental force, with his telepathic wave, that of a wise man. The Master imprinted images which worked on the subject throughout his whole life. These images consisted of visions and trials of life after death, apparitions of demons, monsters, obstacles to overcome, in which took their part fire, earth, air and water, a judgment before the Gods and the final beatific vision.

The Master was able to imprint his will indelibly, and on awakening after the long lethargic sleep,the disciple had acquired a fictitious wisdom and the certainty of having seen the hereafter, of having experienced death and resurrection (spiritual rebirth). With this lifelong suggestion, the disciple remained closely dependent on the master, considered himself twice born or reborn.

After the coming of Christ, spiritual rebirth is only imposed to a minimal extent by the initiator, because except for the influence received as children in baptism from the priest and in Confirmation through the Bishop, when the subtle bodies of the newborn or the child are still impressionable, it is the disciple who must freely seek it with his own personal effort.

The purpose of the earthly master is the awakening of the soul, the spirit and the eros in us in the Light of God, and in all this the Christian has the privilege of being in the atmosphere, in the spiritual aura of the Christ, in his ``Mystical Body'', because the Man Jesus has been penetrated by the Word. There is still more. The Christ is now a force co-present in all humankind, regardless of whether one has been, and we repeat it, has been baptized or not. We are all ``Christophers''
\end{quotationx}


Finally,

\begin{quotex}
... according to the doctrine that the Zendic Magi had given to Jesus during his childhood and to his precursor John the Baptist. The visit and the sojourn of the Magi Kings, Astrologers(1) diviners, therapeutists, is connected with the didactic preparation of Jesus and of the Baptist. The Zendic Magi had learned from the Veda and condensed in the Avesta the secret doctrine. This was learned by an inner circle of Hebrews during the time of the deportation in Babylon, and transmitted orally and partly in writing in the Kabbalistic literature and in the incomplete writings of the Ecclesiasticus, the Book of Wisdom, the Book of Daniel and the Book of Enoch: doctrines that passed to the secret Orders of the Therapeutists and the Essenes, and from them to John the Baptist and the Messiah ... \end{quotex}



\flrightit{Posted on 2011-09-30 by Logres}

\begin{center}* * *\end{center}

\begin{footnotesize}
\begin{sffamily}
\texttt{Charlotte Cowell on 2011-10-02 at 15:51 said:}

Interesting post and helpful reminder about meditation, something I stopped doing a year or so ago (by choice), but have been wishing to return to. Speaking of the royal road, I'm instantly put in mind of Attar's wonderful poem, Bird Parliament, an extract of which is here:

\url{http://alchemical-weddings.com/alchemical-weddings/the-royal-road}


and this is the preceding section of this masterpiece of Sufi literature:

\url{http://alchemical-weddings.com/alchemical-weddings/lions-heart-beneath-the-wing}


Speaking of birds, I'm soon to be traveling to Guatemala so have started to acquaint myself with the marvelous but endangered Quetzal bird... .

\hfill

\texttt{Eagle on 2011-10-02 at 16:52 said:}

Bardon borrows freely from Eastern sources, actually. Representing it as a wholey Western tradition passed down is simply false.

Claiming Christian initation is `superior' to its predecessors is absurd.

Alchemy was not orignally developed by Christians.

So many fallacies in this article.

\hfill

\texttt{Logres on 2011-10-02 at 18:18 said:}

Palamidessi never claims it is ``superior'' -- once again, literal thinking in evidence here. Something can be ``higher'' without in any sense being superior. Other things can perfect or complete what occured before. Is old age superior to the peak of manhood? Which season of the year is better than the others? Palamidessi is claiming that Christ is the completion of what already is. As to Bardon's ``eastern-ness'' that may well be true -- what difference does it make to the point being made? Is he a false or dead end? Does he distort the teaching? Cologero's original point is about the New Age effort to go ``eastern''. Nothing at all was said about the ``adepts'' or the ``masters'' themselves, who obviously work with various streams of influence.

\hfill

\texttt{Logres on 2011-10-02 at 20:06 said:}

I still appreciate the input, and certainly, I suppose I could always phrase things more carefully; re-reading the article, the point is that East -- West conflicts are merely apparent (unless there is actual occult distortion involved). Bardon's ability to state ``Eastern'' truth in terms consonant with alchemy simply underscores this.

\hfill

\texttt{Charlotte Cowell on 2011-10-04 at 10:13 said:}

Regarding what came `first' or `last' or anywhere else in between, bear in mind that the time-space continuum does not exist in the dimension of pure Spirit, certainly not in the way we understand and experience it here. The astral world is similarly void of the time-space restrictions -- will it and you are just `there' (or not, if you willed imperfectly and lacked proper motivation... ). The historical context is useful -- indeed, it's essential -- if we're really going to consolidate our knowledge of such things during our time on Earth -- but it shouldn't bind us. In my opinion the experience of Christian initiation has -- in the most free and sublime sense -- is that it incorporates all other absolute spiritual truths, including some of the essences that flourish under the banners of Sufism, Judaism, Buddhism (for physical preparation especially) and even Shamanism in really quite a fundamental way. One of the only places I ever disagreed with the author of Meditations on the Tarot is in his treatment of shamanism, which I concluded he did not fully understand.

\hfill

\texttt{Charlotte Cowell on 2011-10-04 at 10:14 said:}

I just read back and realised I wrote an incomplete sentence. I was trying to explain why I feel the Christian initiation is pre-eminent, but perhaps it depends on one's experience... .

\hfill

\texttt{Caleb on 2011-10-16 at 17:04 said:}

The implicit endorsement of the quotes that pagan initiatory experience was only an illusion comes across as diminishing of them. I don't think we're in a place to judge the quality of their ancient inner experiences, though I do agree with the general thrust, that Christian esoteric initiation is higher. Thanks for providing more illumination on early hidden currents in the Church.

Do you have a link to the primary source for Palamadessi's quote from St. Basil? It's by far the best `smoking gun' I've seen, so I tried to look it up, but all I could find in Chapter 17 of De Spiritu Sancto was a discourse on enumeration. Thanks for any help you can offer finding that primary source.

\hfill

\texttt{Caleb on 2011-10-16 at 23:58 said:}

Never mind the last question. A Catholic friend helped me out and found the reference was missing an X and is actually from chapter XXVII.

\hfill

\texttt{HOO on 2011-10-17 at 14:46 said:}

--``though I do agree with the general thrust, that Christian esoteric initiation is higher.''

hahaha. why?

\hfill

\texttt{Caleb on 2011-10-17 at 19:02 said:}

Fair question. Short answer: Christ realized on the material realm what initiates had only achieved at the spiritual.

I respect and use paganism, but until it can point to a practitioner whom achieved self-resurrection on the material plane with as high a confidence interval as Jesus, I'll continue to consider the body of Christ a superior manifestation of the divine.

\hfill

\texttt{HOO on 2011-10-18 at 08:55 said:}

Most likely this is a sentimental issue for you but... truth is truth and Jesus' descriptions corresponds to him being one of those dudes:

´The Tantras go on to reaffirm the ancient Vedic-brahmanic and Buddhist belief in the siddha's superiority to all other ``divinities.''

The siddha has power over ``three worlds.'' No god, including Brahma, Vishnu, and Hari-Hara, can resist him. The possibility of doing as one pleases and to prevent things from happening is theoretically upheld. The siddha is ``Lord of Death'' (mrityunjeya) in the specific sense of being able to kill the body through an act of the will -- the so-called samhara-mudra, the ``act of dissolution'' -- of not experiencing death, and of transferring one's consciousness to any chosen level of existence. A siddha, according to Milarepa, ``goes at will through existences as an untamed lion freely roams a mountainous region.''
... In particular, a Siddha cannot lie, since his word is a word of

power, which commands to reality; thus everything he says would come true. In him the fundamental Tantric motif of the unity of bhoga and mukti is actualized. He enjoys the dignity of chakravarti, of a ``World Ruler,''´

Lord of the World, King of Kings, eh? : )

nothing new.

\hfill

\texttt{Boreas on 2011-10-18 at 10:37 said:}

``Lord of the World, King of Kings, eh? : ) nothing new.''

By his own words, he did not came to abolish ``the prophets'' (I'd prefer the word seer, a rohan) but ``to fulfill them'', that is, to proclaim the new covenant and start a new era \& path in spiritual history. This takes nothing out of the seers of the past; it is not a competition about who's the best. The laws and child plays of the mundane world don't apply to the spiritual reality; there is no competition between true masters of the path, only a mutual understanding and assistance.

As a sidenote, talking about these things (whether christian or pagan initiation is better) is especially limited to speculation if oneself hasn't gone through a spiritual rebirth and initiation.

\hfill

\texttt{HOO on 2011-10-18 at 12:26 said:} 

Boreas: ´proclaim the new covenant´ 

what is so new about it? Boreas: ´and start a new era´

yeah, what an era! that was started...  Boreas: ´The laws and child plays of the mundane world don't apply to the spiritual reality´

What are you saying here? What about ``So above, so below''. The world is the playground of the gods, the difference between realms is in subtleness, density, fluidity. When my body is less dense and more conducive to the spirit I move differently and people receive my vibrations differently.

The mundane world is inside the spiritual-reality. Heaven is the Celestial Earth. There is one reality of different levels. Crude nature is a veil in front of Heavenly ``Nature''.

´it is the soul that encloses and bears the body. This is why it is not possible to say where the spiritual place is situated; it is not situated, it is, rather, that which situates, it is situative. Its ubi is an ubique. Certainly, there may be topographical correspondences between the sensory world and the mundus imaginalis, one symbolizing with the other.´ ``Mundus Imaginalis--or the Imaginary and the Imaginal'', Henry Corbin  Boreas: ''As a sidenote, talking about these things (whether christian or pagan initiation is better) is especially limited to speculation if oneself hasn't gone through a spiritual rebirth and initiation.''

Just my point toward Caleb's ``Christian esoteric initiation is higher.'' : )

Notice Logres' last quote in his post corresponds to my previous comment.

\hfill

\texttt{Boreas on 2011-10-18 at 15:38 said:}

``what is so new about it?''

There are quite a few of them, the most fundamental being according to my understanding the almost complete upheaval in the practical methodology and requirements. If I understand correctly, the initiates of the old pagan world had to follow a very precise and regimented approach with many outward practices, the initiations were taken hierarchically one at a time, but the christic initiation and path opened by Christ emphasized the `fourth initiation' (the way of the heart) and made it the most important one of all, since after this is truly traversed -- the rosi-crucian mystery of the bleeding heart -- and the sacrifice made, nothing can make man fall again. The disciple of the new covenant is not necessarily forced to go through a series of lesser initiations before the fourth one, for example in the astral world (which can be even very dangerous in these times), although these paths are of course still open to those who wish to follow the gradual purification of the lower vehicles prior to the fourth initiation. The new path stresses the internal sacrifice and crucifixion more than the former.

There are numerous other issues also, too difficult and lenghty to go here, but I believe this being the most important and fundamental.

``yeah, what an era! that was started... ''

You are here confounding historical christianity and esoteric / mystical side of the christic initiation that are basically without any relevance to each other.

``And behold, I shall be with you until the consummation of the world.''

``My time is not yet, but yours is always.''

``I shall throw fire into the world and watch it until it blazes.''

``What are you saying here?''

Masters of different systems don't compete with each other like politicians in the parlamentary democratic party system, by trying to pull the rug off from their collegues -- even if their wordly representatives and disciples sometimes do this for better or for worse. Instead, their teachings support each other and work for the same goal.


\hfill

\end{sffamily}
\end{footnotesize}

